\begin{abstract}
Interpretability methods produce measurements, but a measurement does not by itself establish the claim commonly attached to it. We present BeyondNN, an evidence and audit layer for PyTorch that records typed, provenance-bound evidence, represents claims about neural networks by relation, subject, scope and eligible units, and re-derives each claim's standing under declared \textsc{primary}, \textsc{alternative} and \textsc{stress\_test} configurations. We give a formal semantics for evidence admissibility, test outcomes and standings, with six properties that hold by construction. On six held-out sites of trained MLP, CNN, BERT-tiny and BERT-base models, 47 of 52 \textsc{primary}-supported Integrated Gradients necessity claims had at least one contradicting declared alternative; an exact re-execution of the frozen design shows that \TotFixed{} of the \TotRev{} reversals hold the selected units fixed, while \TotKOnly{} arise only from changing the selection size and therefore the claim's subject. On compiled Tracr programs the primary test supported 0 of 400 structurally decoy tests and 591 of 760 used-component tests; on 18 trained InterpBench transformers it supported \IbTp{} of \IbPos{} specification-necessary and \IbFp{} of \IbNeg{} specification-unnecessary instances, all \IbFp{} on nodes the trained networks demonstrably use. On identical measurements with seven injected workflow faults, single-configuration reporting surfaced none; a conventional checker found all seven only with a check written for each, and the audit found all seven from its declared plan. The central findings replicated across five independently trained checkpoints per trained model. BeyondNN does not establish a model's mechanism; it makes the assumptions and counterevidence behind a declared test explicit and checkable.
\end{abstract}

\begin{IEEEkeywords}
Interpretability, explainability, faithfulness, causal interventions, attribution, provenance, reproducibility, neural networks.
\end{IEEEkeywords}
